%
% k1.tex -- Kreis k1
%
% (c) 2021 Prof Dr Andreas Müller, OST Ostschweizer Fachhochschule
%
\documentclass[tikz]{standalone}
\usepackage{amsmath}
\usepackage{times}
\usepackage{txfonts}
\usepackage{pgfplots}
\usepackage{csvsimple}
\usetikzlibrary{arrows,intersections,math}
\begin{document}
\def\skala{1}
\begin{tikzpicture}[>=latex,thick,scale=\skala]

% k0: Mittelpunkt M0=(0,0), Radius r0=1.0
% k1: Mittelpunkt M1=(0,1.0), Radius r1=2.0
% Z=(0,-1.0), P1,1=(2.0,1.0), P1,2=(-2.0,1.0)
\def\rz{1.4}  % Radius k0
\pgfmathparse{2*\rz}
\def\ro{\pgfmathresult}  % Radius k1
\coordinate (M0)  at (0,    0);
\coordinate (M1)  at (0,  \rz);
\coordinate (Z)   at (0, -\rz);
\coordinate (P11) at ( \ro, \rz);
\coordinate (P12) at (-\ro, \rz);
% Kreise
\draw (M0) circle (\rz);
\draw (M1) circle (\ro);
% Horizontale Linie durch M1 (Senkrechte zu ZM1)
\draw[thin] (P12) -- (P11);
% Linien von Z zu P1,1 und P1,2
\draw (Z) -- (P11);
\draw (Z) -- (P12);
% Punkte – Labels mit weissem Hintergrund, kollisionsfrei
\fill (M0)  circle (2pt) node[right=6pt, fill=white, inner sep=1pt] {$M_0$};
\fill (M1)  circle (2pt) node[right=6pt, fill=white, inner sep=1pt] {$M_1$};
\fill (Z)   circle (2pt) node[below=4pt, fill=white, inner sep=1pt] {$Z$};
\fill (P11) circle (2pt) node[right=6pt, fill=white, inner sep=1pt] {$P_{1,1}\mathstrut$};
\fill (P12) circle (2pt) node[left=6pt,  fill=white, inner sep=1pt] {$P_{1,2}\mathstrut$};
% k1-Label mit weissem Hintergrund
\node[fill=white, inner sep=1pt] at (-2.6, 3.0) {$k_1$};

\end{tikzpicture}
\end{document}

